\bookchapter{}{References}{The sources of the theorems that this book states without a full proof.}
\label{app:refs}

The lessons cite these works by author and year.

\begin{description}[leftmargin=1.5em, itemindent=-1.5em, itemsep=5pt, font=\normalfont]
  \item[] B. Açıkmeşe, S.~R. Ploen, \enquote{Convex programming approach to powered descent guidance for Mars landing}, \emph{Journal of Guidance, Control, and Dynamics} 30(5), 2007. Rocket landing as a convex optimisation, solved in flight.
  \item[] B.~D.~O. Anderson, J.~B. Moore, \emph{Optimal Control: Linear Quadratic Methods}, Prentice Hall, 1990. The LQR, its Riccati equation (ch.~3) and its robustness (ch.~5).
  \item[] B.~D.~O. Anderson, J.~B. Moore, \emph{Optimal Filtering}, Prentice Hall, 1979. The Kalman filter (ch.~3), its steady state and its duality with the regulator (ch.~4).
  \item[] K.~J. Åström, R.~M. Murray, \emph{Feedback Systems: An Introduction for Scientists and Engineers}, 2nd ed., Princeton University Press, 2021. The closest companion to this book: the same subject with the same respect for both intuition and proof.
  \item[] K.~J. Åström, T. Hägglund, \emph{Advanced PID Control}, ISA, 2006. The practical side of PID in full: anti-windup (ch.~3), setpoint weighting, filtering, tuning rules and their limits.
  \item[] K.~J. Åström, B. Wittenmark, \emph{Adaptive Control}, 2nd ed., Addison-Wesley, 1995. Recursive least squares with forgetting, persistent excitation and the convergence of adaptive controllers (ch.~2).
  \item[] A. Bemporad, M. Morari, V. Dua, E.~N. Pistikopoulos, \enquote{The explicit linear quadratic regulator for constrained systems}, \emph{Automatica} 38(1), 2002. Constrained linear MPC is a piecewise affine state feedback.
  \item[] S.~P. Bhat, D.~S. Bernstein, \enquote{A topological obstruction to continuous global stabilization of rotational motion and the unwinding phenomenon}, \emph{Systems \& Control Letters} 39(1), 2000. Why no continuous attitude law can bring every attitude to rest.
  \item[] L. Blackmore, \enquote{Autonomous precision landing of space rockets}, \emph{The Bridge} 46(4), National Academy of Engineering, 2016. How reusable boosters land, by the engineer who led the work.
  \item[] H.~W. Bode, \emph{Network Analysis and Feedback Amplifier Design}, Van Nostrand, 1945. Frequency-response design, the margins and the sensitivity integral.
  \item[] S. Boyd, L. Vandenberghe, \emph{Convex Optimization}, Cambridge University Press, 2004. Convex sets and functions, duality, the KKT conditions, cone programs.
  \item[] D. Brescianini, R. D'Andrea, \enquote{Tilt-prioritized quadrocopter attitude control}, \emph{IEEE Transactions on Control Systems Technology} 28(2), 2020. The split of the attitude error into tilt and yaw.
  \item[] G. Cybenko, \enquote{Approximation by superpositions of a sigmoidal function}, \emph{Mathematics of Control, Signals, and Systems} 2(4), 1989. One hidden layer suffices to approximate any continuous function on a compact set.
  \item[] J.~C. Doyle, \enquote{Guaranteed margins for LQG regulators}, \emph{IEEE Transactions on Automatic Control} 23(4), 1978. The abstract reads, in full: \enquote{There are none.}
  \item[] J.~C. Doyle, B.~A. Francis, A.~R. Tannenbaum, \emph{Feedback Control Theory}, Macmillan, 1992. The sensitivity functions and the limits on what feedback can do, including Bode's integral (ch.~6).
  \item[] W.~R. Evans, \enquote{Graphical analysis of control systems}, \emph{Transactions of the AIEE} 67(1), 1948. The root locus.
  \item[] M. Faessler, A. Franchi, D. Scaramuzza, \enquote{Differential flatness of quadrotor dynamics subject to rotor drag for accurate tracking of high-speed trajectories}, \emph{IEEE Robotics and Automation Letters} 3(2), 2018. Rotor drag, proportional to thrust and to the in-plane velocity.
  \item[] T. Flash, N. Hogan, \enquote{The coordination of arm movements: an experimentally confirmed mathematical model}, \emph{Journal of Neuroscience} 5(7), 1985. Human reaching follows the minimum-jerk path.
  \item[] M. Fliess, J. Lévine, P. Martin, P. Rouchon, \enquote{Flatness and defect of non-linear systems: introductory theory and examples}, \emph{International Journal of Control} 61(6), 1995. Differential flatness.
  \item[] B.~A. Francis, W.~M. Wonham, \enquote{The internal model principle of control theory}, \emph{Automatica} 12(5), 1976. Why a loop must contain a model of what it is to reject.
  \item[] G.~F. Franklin, J.~D. Powell, A. Emami-Naeini, \emph{Feedback Control of Dynamic Systems}, 8th ed., Pearson, 2019. The standard undergraduate text; root locus (ch.~5), frequency response (ch.~6).
  \item[] F.~R. Gantmacher, \emph{The Theory of Matrices}, vol.~2, Chelsea, 1959. The Routh–Hurwitz problem in full (ch.~XV).
  \item[] G.~H. Golub, C.~F. Van Loan, \emph{Matrix Computations}, 4th ed., Johns Hopkins University Press, 2013. The SVD, Cholesky, and how matrix equations are solved in practice.
  \item[] E. Hairer, S.~P. Nørsett, G. Wanner, \emph{Solving Ordinary Differential Equations I: Nonstiff Problems}, 2nd ed., Springer, 1993. Order, convergence and stability of one-step methods.
  \item[] M.~W. Hirsch, S. Smale, R.~L. Devaney, \emph{Differential Equations, Dynamical Systems, and an Introduction to Chaos}, 3rd ed., Academic Press, 2013. Linear systems, the matrix exponential and phase portraits, with proofs.
  \item[] R.~A. Horn, C.~R. Johnson, \emph{Matrix Analysis}, 2nd ed., Cambridge University Press, 2013. Symmetric and positive definite matrices, the spectral theorem.
  \item[] K. Hornik, M. Stinchcombe, H. White, \enquote{Multilayer feedforward networks are universal approximators}, \emph{Neural Networks} 2(5), 1989. Universal approximation for a wide class of activation functions.
  \item[] T. Kailath, \emph{Linear Systems}, Prentice-Hall, 1980. State space, controllability and pole placement in full.
  \item[] R.~E. Kálmán, \enquote{A new approach to linear filtering and prediction problems}, \emph{Journal of Basic Engineering} 82(1), 1960. The Kalman filter.
  \item[] R.~E. Kálmán, \enquote{On the general theory of control systems}, \emph{Proceedings of the First IFAC Congress}, Moscow, 1960. Controllability and observability.
  \item[] R.~E. Kálmán, \enquote{When is a linear control system optimal?}, \emph{Journal of Basic Engineering} 86(1), 1964. The return-difference inequality behind the margins of the LQR.
  \item[] E. Kaufmann, L. Bauersfeld, A. Loquercio, M. Müller, V. Koltun, D. Scaramuzza, \enquote{Champion-level drone racing using deep reinforcement learning}, \emph{Nature} 620, 2023. A learned policy, commanding thrust and body rates, beats human champions.
  \item[] H.~K. Khalil, \emph{Nonlinear Systems}, 3rd ed., Prentice Hall, 2002. Existence and uniqueness (ch.~3), Lyapunov stability, linearisation and the invariance principle (ch.~4), averaging (ch.~10), singular perturbations (ch.~11).
  \item[] P. Kunapuli, J. Welde, D. Jayaraman, V. Kumar, \enquote{Leveling the playing field: carefully comparing classical and learned controllers for quadrotor trajectory tracking}, \emph{Robotics: Science and Systems}, 2025 (arXiv:2506.17832). Learned policies win transients, geometric control the steady state.
  \item[] A.~J. Laub, \enquote{A Schur method for solving algebraic Riccati equations}, \emph{IEEE Transactions on Automatic Control} 24(6), 1979. The Hamiltonian method as numerical software uses it.
  \item[] T. Lee, M. Leok, N.~H. McClamroch, \enquote{Geometric tracking control of a quadrotor UAV on SE(3)}, \emph{49th IEEE Conference on Decision and Control}, 2010. Attitude and position control written on the rotation group, with almost-global stability.
  \item[] L. Ljung, T. Söderström, \emph{Theory and Practice of Recursive Identification}, MIT Press, 1983. Recursive least squares and its relatives, with their convergence analysis.
  \item[] F.~L. Markley, J.~L. Crassidis, \emph{Fundamentals of Spacecraft Attitude Determination and Control}, Springer, 2014. Rotation matrices, quaternions and their kinematics (ch.~2, 3), with every convention stated.
  \item[] D.~Q. Mayne, J.~B. Rawlings, C.~V. Rao, P.~O.~M. Scokaert, \enquote{Constrained model predictive control: stability and optimality}, \emph{Automatica} 36(6), 2000. Recursive feasibility and stability of MPC with terminal ingredients.
  \item[] D. Mellinger, V. Kumar, \enquote{Minimum snap trajectory generation and control for quadrotors}, \emph{IEEE International Conference on Robotics and Automation}, 2011. The quadrotor is flat; trajectories as polynomials of minimum snap.
  \item[] J. Nocedal, S.~J. Wright, \emph{Numerical Optimization}, 2nd ed., Springer, 2006. Optimality conditions, gradient and Newton methods, quadratic programming.
  \item[] H. Nyquist, \enquote{Regeneration theory}, \emph{Bell System Technical Journal} 11(1), 1932. The stability criterion that bears his name.
  \item[] M. O'Connell, G. Shi, X. Shi, K. Azizzadenesheli, A. Anandkumar, Y. Yue, S.-J. Chung, \enquote{Neural-Fly enables rapid learning for agile flight in strong winds}, \emph{Science Robotics} 7(66), 2022. Learned features of the aerodynamic residual, linear coefficients adapted online.
  \item[] A.~V. Oppenheim, A.~S. Willsky, \emph{Signals and Systems}, 2nd ed., Prentice Hall, 1997. Fourier series and transforms, the Laplace and z-transforms, with all the proofs this book omits.
  \item[] A.~V. Oppenheim, R.~W. Schafer, \emph{Discrete-Time Signal Processing}, 3rd ed., Pearson, 2010. Sampling, aliasing and digital filters.
  \item[] A. Papoulis, S.~U. Pillai, \emph{Probability, Random Variables and Stochastic Processes}, 4th ed., McGraw-Hill, 2002. Random processes, spectra and linear systems (ch.~9), estimation (ch.~8).
  \item[] D.~A. Pomerleau, \enquote{ALVINN: an autonomous land vehicle in a neural network}, \emph{Advances in Neural Information Processing Systems} 1, 1989. A neural network steers a vehicle from camera images.
  \item[] D.~A. Pomerleau, \enquote{Efficient training of artificial neural networks for autonomous navigation}, \emph{Neural Computation} 3(1), 1991. Learning to steer by watching a driver, with synthesised views to learn recovery.
  \item[] S.~J. Qin, T.~A. Badgwell, \enquote{A survey of industrial model predictive control technology}, \emph{Control Engineering Practice} 11(7), 2003. MPC in the process industries, from the 1970s on.
  \item[] A. Quarteroni, R. Sacco, F. Saleri, \emph{Numerical Mathematics}, 2nd ed., Springer, 2007. Root finding, quadrature and linear algebra with their error analysis.
  \item[] Lord Rayleigh (J.~W. Strutt), \enquote{The soaring of birds}, \emph{Nature} 27, 1883, pp.~534--535; reprinted in his \emph{Scientific Papers}, vol.~II, Cambridge University Press, 1900. Soaring on the change of the wind with height: dynamic soaring.
  \item[] S. Ross, J.~A. Bagnell, \enquote{Efficient reductions for imitation learning}, \emph{Proceedings of the 13th International Conference on Artificial Intelligence and Statistics (AISTATS)}, 2010. Why errors of behaviour cloning compound quadratically with the horizon.
  \item[] S. Ross, G.~J. Gordon, J.~A. Bagnell, \enquote{A reduction of imitation learning and structured prediction to no-regret online learning}, \emph{Proceedings of the 14th International Conference on Artificial Intelligence and Statistics (AISTATS)}, 2011. DAgger: train on the states the learner itself visits.
  \item[] M.~G. Safonov, M. Athans, \enquote{Gain and phase margin for multiloop LQG regulators}, \emph{IEEE Transactions on Automatic Control} 22(2), 1977. The margins of the LQR for several inputs.
  \item[] T. Salimans, J. Ho, X. Chen, S. Sidor, I. Sutskever, \enquote{Evolution strategies as a scalable alternative to reinforcement learning}, arXiv:1703.03864, 2017. Antithetic Gaussian perturbations and rank-based fitness shaping for training policies.
  \item[] C.~E. Shannon, \enquote{Communication in the presence of noise}, \emph{Proceedings of the IRE} 37(1), 1949. The sampling theorem.
  \item[] E.~J.~J. Smeur, Q.~P. Chu, G.~C.~H.~E. de Croon, \enquote{Adaptive incremental nonlinear dynamic inversion for attitude control of micro air vehicles}, \emph{Journal of Guidance, Control, and Dynamics} 39(3), 2016. INDI on a quadrotor, with the synchronised filters and the stability analysis of the increment loop.
  \item[] B. Stellato, G. Banjac, P. Goulart, A. Bemporad, S. Boyd, \enquote{OSQP: an operator splitting solver for quadratic programs}, \emph{Mathematical Programming Computation} 12(4), 2020. The ADMM solver behind the simulator's MPC.
  \item[] S. Sun, A. Romero, P. Foehn, E. Kaufmann, D. Scaramuzza, \enquote{A comparative study of nonlinear MPC and differential-flatness-based control for quadrotor agile flight}, \emph{IEEE Transactions on Robotics} 38(6), 2022. Each wins a different regime; an INDI inner loop makes both robust.
  \item[] S. Thrun, M. Montemerlo, H. Dahlkamp, D. Stavens, et al., \enquote{Stanley: the robot that won the DARPA Grand Challenge}, \emph{Journal of Field Robotics} 23(9), 2006. Learned perception and classical control in one vehicle.
  \item[] H. Weimerskirch, K. Delord, A. Guitteaud, R.~A. Phillips, P. Pinet, \enquote{Extreme variation in migration strategies between and within wandering albatross populations during their sabbatical year, and their fitness consequences}, \emph{Scientific Reports} 5, 8853, 2015. Albatrosses that circle Antarctica two or three times in a year.
  \item[] W.~M. Wonham, \enquote{On pole assignment in multi-input controllable linear systems}, \emph{IEEE Transactions on Automatic Control} 12(6), 1967. Poles can be placed anywhere exactly when the system is controllable.
  \item[] L. Zaccarian, A.~R. Teel, \emph{Modern Anti-windup Synthesis}, Princeton University Press, 2011. Stability guarantees for loops with saturating actuators.
\end{description}
